\chapter{Rasterização}

A rasterização é o processo fundamental de conversão de descrições geométricas contínuas em uma grade discreta de pixels, permitindo a exibição de cenas 3D em monitores digitais. Após as etapas de transformações de visualização e projeção, os vértices que compõem os objetos são mapeados para o \textit{viewport} em coordenadas de tela. Contudo, como o hardware de exibição moderno é composto por uma matriz finita de elementos de cor, denominados pixels, é necessário realizar uma amostragem sistemática das primitivas geométricas. Como destaca \textcite{shirley2021fundamentals}, a rasterização deve equilibrar a fidelidade visual com uma eficiência computacional extrema, processando bilhões de fragmentos por segundo em arquiteturas de GPU altamente paralelizadas.

\section{O Problema da Discretização}

O hardware gráfico opera em um espaço de coordenadas de tela, onde cada pixel é uma amostra pontual em uma grade regular. O problema central da rasterização é determinar quais pixels são "cobertos" por uma primitiva geométrica. Este processo envolve não apenas a geometria, mas também a interpolação de atributos (como cores e texturas) e a resolução da visibilidade (quais objetos estão na frente).

\section{Rasterização de Pontos e Linhas}

Embora triângulos sejam as primitivas dominantes, a rasterização de pontos e linhas é a base para ferramentas de CAD e debug visual.

\subsection{Rasterização de Pontos}

Um ponto matemático não possui área. No entanto, para ser visível, ele deve ser mapeado para pelo menos um pixel. A regra de rasterização de pontos geralmente ativa o pixel $(x, y)$ tal que $x = \lfloor x_p \rfloor$ e $y = \lfloor y_p \rfloor$. Em aplicações onde pontos representam partículas, pode-se utilizar um \textit{Point Size} maior que 1.0, o que resulta na ativação de um bloco de pixels ao redor da coordenada central.

\subsection{O Algoritmo de Bresenham para Linhas}

O traçado de linhas em uma grade discreta apresenta o desafio de escolher os pixels que melhor aproximam uma trajetória contínua. O algoritmo de Bresenham, desenvolvido em 1962, é a solução clássica que evita operações de ponto flutuante e divisões, utilizando apenas incrementos inteiros.

\subsubsection{Derivação Matemática e a Variável de Decisão}

Consideremos uma linha no primeiro octante, onde a inclinação $m = \Delta y / \Delta x$ satisfaz $0 \le m \le 1$. Ao avançar uma unidade no eixo $X$, devemos decidir se o pixel no eixo $Y$ permanece o mesmo ou incrementa para $y+1$.

Definimos o erro acumulado $e$ como a distância vertical entre a posição real da linha e a posição discreta do pixel. Para otimização, trabalhamos com uma variável de decisão $d$ que elimina frações. Ao analisar o ponto médio entre os pixels candidatos, derivamos:
\begin{itemize}
    \item Se $P_i \le 0$: Escolhemos o pixel inferior (mesmo $Y$). $P_{i+1} = P_i + 2\Delta y$.
    \item Se $P_i > 0$: Escolhemos o pixel superior ($Y+1$). $P_{i+1} = P_i + 2(\Delta y - \Delta x)$.
\end{itemize}

O valor inicial é $P_0 = 2\Delta y - \Delta x$. Esta abordagem garante que a linha rasterizada seja a representação mais fiel possível dentro da grade, minimizando o erro de aproximação visual (\textit{aliasing} geométrico).

\subsubsection{Generalização para Todos os Octantes}

Para suportar todas as orientações possíveis, aplicamos transformações de simetria:
1. Se $|\Delta y| > |\Delta x|$, trocamos os eixos $X$ e $Y$ (linha predominantemente vertical).
2. Se $x_0 > x_1$, invertemos a ordem dos pontos de controle.
3. Se $y_0 > y_1$, o incremento em $Y$ torna-se negativo ($step_y = -1$).

\begin{lstlisting}[language=C++, caption=Algoritmo de Bresenham Robusto]
void drawLine(int x0, int y0, int x1, int y1) {
    int dx = abs(x1 - x0), sx = x0 < x1 ? 1 : -1;
    int dy = -abs(y1 - y0), sy = y0 < y1 ? 1 : -1;
    int err = dx + dy;

    while (true) {
        setPixel(x0, y0, color);
        if (x0 == x1 && y0 == y1) break;
        int e2 = 2 * err;
        if (e2 >= dy) { err += dy; x0 += sx; }
        if (e2 <= dx) { err += dx; y0 += sy; }
    }
}
\end{lstlisting}

\section{Rasterização de Triângulos}

Triângulos são as primitivas fundamentais das GPUs modernas por serem garantidamente planos e convexos.

\subsection{Bounding Box e Teste de Inclusão}

O pipeline moderno de rasterização utiliza uma abordagem de força bruta inteligente. Para cada triângulo:
1. Calcula-se a \textbf{Axis-Aligned Bounding Box (AABB)} no espaço de tela.
2. Clampa-se a caixa contra os limites do \textit{viewport}.
3. Para cada pixel $(x, y)$ na caixa, aplica-se a \textbf{Função de Aresta} ($E$):
\begin{equation}
    E(P) = (x - x_i)(y_j - y_i) - (y - y_i)(x_j - x_i)
\end{equation}
Se $E(P) > 0$ para as três arestas, o pixel é coberto pelo triângulo.

\subsection{A Regra de Preenchimento Top-Left}

Para evitar a rasterização dupla de pixels sobre arestas compartilhadas (o que causaria problemas em transparências), as GPUs utilizam a regra \textit{Top-Left}. Um pixel na aresta só pertence ao triângulo se a aresta for uma "aresta de topo" (horizontal e acima) ou uma "aresta esquerda" (situada à esquerda).

\section{Coordenadas Baricêntricas}

Coordenadas baricêntricas $(\lambda_0, \lambda_1, \lambda_2)$ fornecem uma forma natural de descrever a posição de um ponto $P$ em relação aos vértices $V_0, V_1, V_2$.

\subsection{Definição via Áreas com Sinal}

Cada $\lambda_i$ é a razão entre a área do sub-triângulo oposto ao vértice $V_i$ e a área total $A_{ABC}$:
\begin{formula}{Áreas Baricêntricas}
    \lambda_0 = \frac{\text{Area}(P, B, C)}{\text{Area}(A, B, C)}, \quad \lambda_1 = \frac{\text{Area}(A, P, C)}{\text{Area}(A, B, C)}, \quad \lambda_2 = \frac{\text{Area}(A, B, P)}{\text{Area}(A, B, C)}
\end{formula}

Propriedades fundamentais:
\begin{itemize}
    \item $\lambda_0 + \lambda_1 + \lambda_2 = 1$.
    \item $P$ está dentro se $\lambda_i \ge 0$ para todo $i$.
    \item Se um $\lambda_i = 0$, $P$ está sobre a aresta oposta a $V_i$.
\end{itemize}

\subsection{Interpolação de Atributos}

As coordenadas baricêntricas são usadas para interpolar atributos de vértice (cor, normal, UV) de forma suave:
\begin{equation}
    \Phi_P = \lambda_0 \Phi_0 + \lambda_1 \Phi_1 + \lambda_2 \Phi_2
\end{equation}

\section{Visibilidade e o Z-Buffer}

O problema de quais triângulos são visíveis em cada pixel é resolvido pelo \textbf{Z-Buffer} (ou \textit{Depth Buffer}).

\subsection{O Algoritmo do Depth Buffer}

Diferente do Algoritmo do Pintor (que exige ordenação $O(N \log N)$), o Z-Buffer opera por fragmento:
1. Inicializar buffer com $1.0$ (máxima profundidade).
2. Para cada fragmento gerado com profundidade $z \in [0, 1]$:
   - Se $z < Z\_Buffer[x, y]$:
     - $Z\_Buffer[x, y] = z$.
     - Atualizar cor do pixel.

\subsection{Não-Linearidade de Z e Z-Fighting}

A projeção perspectiva mapeia $Z$ de forma não-linear ($1/z$). Isso concentra a precisão perto do plano \textit{near}. O \textbf{Z-fighting} ocorre quando a precisão finita do buffer (ex: 24 bits) impede a distinção entre dois fragmentos muito próximos. Afastar o plano \textit{near} da câmera é a solução primária.

\subsection{Interpolação Perspectiva-Correta}

Interpolar atributos linearmente na tela produz distorções (ex: texturas "balançando"). A solução matemática é interpolar $1/w$ e os atributos divididos por $w$:
\begin{formula}{Correção Perspectiva}
    \Phi_{fragmento} = \frac{\text{linear}(\Phi/w)}{\text{linear}(1/w)}
\end{formula}

\section{Antialiasing}

O \textit{aliasing} (serrilhado) ocorre devido à amostragem insuficiente de sinais de alta frequência (bordas).

\begin{itemize}
    \item \textbf{SSAA (Super-Sampling):} Renderiza em $4\times$ resolução e reduz. Qualidade máxima, custo altíssimo.
    \item \textbf{MSAA (Multi-Sampling):} Otimiza o SSAA executando o fragment shader apenas uma vez por pixel, mas testando cobertura em sub-pixels.
    \item \textbf{FXAA (Fast Approximate):} Filtro de pós-processamento baseado em detecção de bordas por contraste.
    \item \textbf{TAA (Temporal):} Acumula amostras de frames anteriores usando reprojeção com vetores de movimento.
\end{itemize}

\section{Alpha Blending}

A transparência é simulada via composição \textit{Over}:
\begin{formula}{Equação de Blending}
    C_{out} = \alpha_s C_s + (1 - \alpha_s) C_d
\end{formula}
Objetos transparentes devem ser renderizados de trás para frente (\textit{back-to-front}) após todos os opacos. O uso de \textit{Pre-multiplied Alpha} evita artefatos de cor em filtros de textura.

\section*{Exercícios}

\begin{enumerate}
    \item Execute o algoritmo de Bresenham manualmente para a linha $(0, 0)$ a $(5, 3)$.
    \item Deduza a variável de decisão inicial $P_0$ para Bresenham.
    \item Prove que as coordenadas baricêntricas somam 1 para qualquer ponto no plano.
    \item Por que a interpolação linear de atributos na tela falha em projeções perspectivas?
    \item Explique como o Z-Buffer lida com a sobreposição de objetos sem necessidade de ordenação prévia.
    \item O que define a precisão do Z-Buffer e como o Z-fighting se manifesta?
    \item Explique o funcionamento do MSAA e por que ele é preferível ao SSAA.
    \item Como o \textit{Pre-multiplied Alpha} altera a matemática do Alpha Blending?
    \item Discuta o impacto do \textit{Early-Z} no desempenho de fragment shaders complexos.
    \item Dada uma aresta compartilhada, como a regra Top-Left garante que cada pixel pertença a apenas um triângulo?
\end{enumerate}


\begin{lstlisting}[language=C++, caption=Implementação Generalizada de Bresenham]
void rasterizeLine(int x0, int y0, int x1, int y1) {
    int dx = abs(x1 - x0), sx = x0 < x1 ? 1 : -1;
    int dy = -abs(y1 - y0), sy = y0 < y1 ? 1 : -1;
    int err = dx + dy;

    while (true) {
        setPixel(x0, y0, color);
        if (x0 == x1 && y0 == y1) break;
        int e2 = 2 * err;
        if (e2 >= dy) { // Erro em X
            err += dy;
            x0 += sx;
        }
        if (e2 <= dx) { // Erro em Y
            err += dx;
            y0 += sy;
        }
    }
}
\end{lstlisting}

\section{Rasterização de Triângulos}

Triângulos são as primitivas de eleição em computação gráfica por serem inerentemente planos, convexos e bem definidos por apenas três pontos. A rasterização moderna abandonou o método de \textit{scanline} horizontal em favor de algoritmos baseados em funções de aresta, que são mais amigáveis ao processamento paralelo massivo das GPUs.

\subsection{O Algoritmo de Bounding Box}

O processo inicia-se com o cálculo da \textit{Axis-Aligned Bounding Box} (AABB) do triângulo no espaço de tela. Esta caixa delimita a área máxima que o triângulo pode ocupar:
\begin{align}
    x_{min} = \lfloor \min(x_0, x_1, x_2) \rfloor, \quad x_{max} = \lceil \max(x_0, x_1, x_2) \rceil \\
    y_{min} = \lfloor \min(y_0, y_1, y_2) \rfloor, \quad y_{max} = \lceil \max(y_0, y_1, y_2) \rceil
\end{align}

Para cada pixel contido nesta região, executamos uma \textbf{função de aresta} $E(x, y)$. Se o valor da função for positivo para as três arestas do triângulo (considerando uma ordem de enrolamento consistente, como sentido anti-horário), o pixel está dentro da primitiva.

\subsection{Regras de Desempate (Top-Left Rule)}

Um problema comum ocorre em pixels localizados exatamente sobre a aresta compartilhada entre dois triângulos adjacentes. Se ambos os triângulos rasterizarem o pixel, haverá sobreposição desnecessária e artefatos em transparências. Para evitar isso, as GPUs aplicam a "Regra de Topo-Esquerda": um pixel na aresta só é desenhado se a aresta for uma "aresta de topo" (horizontal e acima) ou uma "aresta esquerda" (não horizontal e situada à esquerda do corpo do triângulo).

\section{Coordenadas Baricêntricas}

As coordenadas baricêntricas fornecem um sistema de coordenadas local para o triângulo, permitindo expressar qualquer ponto $P$ como uma combinação linear dos vértices $V_0, V_1$ e $V_2$.

\subsection{Definição Matemática e Áreas}

Um ponto $P$ dentro do triângulo satisfaz a equação:
\begin{formula}{Representação Baricêntrica}
    P = \lambda_0 V_0 + \lambda_1 V_1 + \lambda_2 V_2 \\
    \lambda_0 + \lambda_1 + \lambda_2 = 1, \quad \lambda_i \in [0, 1]
\end{formula}

Cada coordenada $\lambda_i$ pode ser interpretada como a razão entre a área do sub-triângulo oposto ao vértice $i$ e a área total do triângulo original. Utilizando o produto vetorial 2D (determinante), calculamos as áreas com sinal:
\begin{equation}
    \text{Area}(A, B, C) = \frac{1}{2} \det \begin{bmatrix} x_B - x_A & x_C - x_A \\ y_B - y_A & y_C - y_A \end{bmatrix}
\end{equation}

A linearidade dessas funções de área permite que as GPUs calculem as coordenadas baricêntricas de forma incremental conforme percorrem os pixels, otimizando drasticamente o teste de inclusão e a interpolação.

\subsection{Interpolação de Atributos e Sombreamento}

Qualquer atributo associado aos vértices (cor, normal, coordenada de textura) é interpolado para o fragmento $P$ usando as mesmas coordenadas baricêntricas:
\begin{equation}
    \Phi_P = \lambda_0 \Phi_0 + \lambda_1 \Phi_1 + \lambda_2 \Phi_2
\end{equation}

\begin{lstlisting}[language=GLSL, caption=Exemplo de Atributos Interpolados]
// No Fragment Shader, os 'in' sao interpolados baricentricamente
in vec3 vNormal; // Normal interpolada
in vec2 vUV;     // Coordenada de textura interpolada

void main() {
    vec3 n = normalize(vNormal);
    vec4 texColor = texture(uTexture, vUV);
    // ... calculo de iluminacao ...
}
\end{lstlisting}

\section{Visibilidade e o Z-Buffer}

A determinação de quais superfícies são visíveis e quais estão ocultas é resolvida pelo algoritmo de \textit{Z-Buffer} (ou \textit{Depth Buffer}), que superou historicamente o Algoritmo do Pintor.

\subsection{Falhas do Algoritmo do Pintor}

Ordenar triângulos do fundo para a frente (\textit{back-to-front}) é ineficiente e falha em casos de:
\begin{itemize}
    \item Interseções cíclicas (A cobre B, B cobre C, C cobre A).
    \item Triângulos que se trespassam mutuamente.
    \item Custo computacional de ordenação $O(N \log N)$ para milhões de polígonos.
\end{itemize}

\subsection{O Algoritmo Depth-Buffer}

O Z-Buffer mantém um buffer de memória auxiliar com a mesma resolução do buffer de cor. Cada entrada armazena a menor profundidade $z$ encontrada até o momento para aquele pixel.

Procedimento:
\begin{enumerate}
    \item Inicializar o Z-Buffer com o valor máximo ($1.0$ em NDC).
    \item Para cada fragmento gerado na posição $(x, y)$ com profundidade $z$:
    \item Se $z < Z\_Buffer[x, y]$:
        \item $Z\_Buffer[x, y] = z$.
        \item Atualizar o pixel no Color Buffer com a cor do fragmento.
    \item Caso contrário, descartar o fragmento.
\end{enumerate}

\subsection{Interpolação Perspectiva-Correta}

Um detalhe matemático fundamental é que a projeção perspectiva não preserva a linearidade da profundidade. Interpolar atributos linearmente no espaço de tela (espaço 2D após a divisão por $W$) causa distorções graves em texturas, conhecidas como "afunilamento".

A correção exige a interpolação linear do inverso da profundidade $1/w$ e dos atributos modificados $\Phi/w$. O valor final no fragmento é recuperado por:
\begin{formula}{Fórmula da Interpolação Perspectiva}
    \Phi_{fragmento} = \frac{\text{interp}(\Phi/w)}{\text{interp}(1/w)}
\end{formula}

\subsection{Z-Fighting e Early-Z}

O \textbf{Z-fighting} manifesta-se como cintilações em superfícies muito próximas, onde a imprecisão numérica do buffer impede a decisão correta de visibilidade. Como a precisão do Z-Buffer é maior perto do plano \textit{near}, uma solução comum é aumentar o valor do \textit{near plane}.

O \textbf{Early-Z} é uma técnica onde o hardware realiza o teste de profundidade \textbf{antes} de executar o Fragment Shader. Se o teste falhar, o shader nem sequer é iniciado, economizando ciclos preciosos de GPU.

\section{Antialiasing}

Devido à natureza discreta dos pixels, bordas de alta frequência geométrica geram serrilhados (\textit{aliasing}). O Teorema de Nyquist afirma que precisamos de uma taxa de amostragem pelo menos duas vezes superior à frequência máxima do sinal.

\subsection{Técnicas de Suavização}

\begin{itemize}
    \item \textbf{SSAA (Super-Sampling):} Renderiza em alta resolução interna e reduz para o tamanho final. Qualidade perfeita, mas custo de memória e processamento proibitivo ($4\times$ ou mais).
    \item \textbf{MSAA (Multi-Sampling):} Otimização onde o shader de fragmento roda apenas uma vez por pixel, mas os testes de cobertura e profundidade ocorrem em múltiplas sub-amostras. Oferece excelente custo-benefício para bordas.
    \item \textbf{FXAA (Fast Approximate):} Algoritmo de pós-processamento que suaviza bordas baseado em contraste. Muito rápido, mas pode "borrar" texturas finas.
    \item \textbf{TAA (Temporal):} Combina informações de múltiplos quadros no tempo, aplicando um pequeno \textit{jitter} na projeção a cada frame. Atualmente é o padrão na indústria por lidar bem com aliasing de movimento e especularidade.
\end{itemize}

\section{Alpha Blending e Transparência}

Para objetos translúcidos, a cor final é uma composição entre a fonte (\textit{source}) e o que já está no buffer (\textit{destination}). A equação "Over" é a mais comum:
\begin{formula}{Composição Alpha}
    C_{out} = \alpha_s C_s + (1 - \alpha_s) C_d
\end{formula}

Diferente de objetos opacos, a ordem de renderização é vital: objetos transparentes devem ser desenhados \textbf{após} os opacos e ordenados de trás para frente. Uma técnica para evitar artefatos de cor em filtragens é o \textbf{Pre-multiplied Alpha}, onde os canais RGB já são armazenados multiplicados pelo alpha ($C' = \alpha C$).

\section*{Exercícios}

\begin{enumerate}
    \item Execute o algoritmo de Bresenham manualmente para a linha entre $(0, 0)$ e $(5, 3)$. Liste os valores da variável de decisão e os pixels escolhidos.
    \item Em um triângulo com vértices $P_0(0,0)$, $P_1(10,0)$ e $P_2(0,10)$, determine as coordenadas baricêntricas do ponto $P(2, 2)$. Ele está dentro do triângulo?
    \item Deduza a expressão da variável de decisão inicial $d_0$ para o algoritmo de Bresenham.
    \item Por que a interpolação perspectiva-correta é necessária? Explique a diferença entre interpolar no espaço 3D e no espaço de tela.
    \item O que é o \textit{Z-fighting} e como a distribuição não-linear de valores no Z-Buffer contribui para este fenômeno?
    \item Explique o funcionamento do MSAA e por que ele é mais eficiente que o SSAA em termos de execução de fragment shaders.
    \item Como o \textit{Pre-multiplied Alpha} simplifica a equação de composição e por que ele é preferido em pipelines de composição profissional?
    \item Discuta o impacto do \textit{Early-Z} na performance de uma cena com alta \textit{overdraw} (muitos objetos sobrepostos).
\end{enumerate}
